\section{\agnews{} with 20\% Symmetric Label Noise}
\label{sec:results-agnews-noisy}

\input{tables/nlp_agnews_noisy}

\Cref{tab:agnews-noise20} reports the \agnews{} results under 20\% symmetric training-label noise, with validation labels kept clean. This setting produces a clearer separation between peak and final validation accuracy than clean \agnews{}. Flat \adamw{} reaches $92.72\%$ peak accuracy but falls to $90.12\%$ final accuracy, corresponding to a $2.59$ percentage-point drop. \adamw{} with warmup-linear scheduling improves the peak to $93.19\%$, but still falls to $90.79\%$ final accuracy, with a $2.40$ percentage-point drop.

In contrast, \method{} substantially improves final validation accuracy. Without an external scheduler, \method{} reaches $93.32\%$ peak accuracy and $92.85\%$ final accuracy, reducing the best-to-final drop to $0.47$ percentage points. When combined with warmup-linear scheduling, \method{} reaches the strongest result in the main table, with $93.55\%$ peak accuracy and $93.51\%$ final accuracy. The best-to-final drop is reduced to only $0.05$ percentage points. This indicates that \method{} does not merely improve the best checkpoint in this setting, but also stabilizes performance toward the end of training.

The axis ablation in \Cref{tab:agnews-noise20-axis} shows that this improvement is primarily driven by adaptive gradient-noise control. \method{}-LR with warmup-linear scheduling reaches only $92.72\%$ peak accuracy and $91.53\%$ final accuracy, which is below the full \method{} variant and also below \method{}-Noise. In contrast, \method{}-Noise reaches $93.47\%$ peak accuracy and $93.45\%$ final accuracy, nearly matching \method{}-Dual at $93.55\%$ peak and $93.51\%$ final accuracy. The best-to-final drop is also almost eliminated for both noise-enabled variants, with $0.01$ percentage points for \method{}-Noise and $0.05$ percentage points for \method{}-Dual.

These results make noisy \agnews{} the strongest NLP evidence for \method{}. Unlike clean \agnews{}, where warmup-linear scheduling already captures most of the available gain, the noisy-label setting leaves substantial room for adaptive control. The results suggest that adaptive gradient-noise control can improve late-training stability under noisy supervision, while adaptive learning-rate modulation alone is not sufficient in this setting.

\begin{figure}[t]
    \centering
    \thesisfigure{0.92\textwidth}{figures/nlp/agnews_noisy_curves.pdf}
    \caption[
        Validation trajectories on \agnews{} with 20\% symmetric label noise
    ]{
        Validation trajectories on \agnews{} with 20\% symmetric training-label noise. \method{} variants with adaptive gradient-noise control reduce late-stage degradation compared with \adamw{} and \adamw{} with warmup-linear scheduling, while the learning-rate-only variant is less effective.
    }
    \label{fig:agnews-noisy-curves}
\end{figure}